\documentclass[12pt]{article}

\usepackage{polyglossia}
\setdefaultlanguage[numerals=thai]{thai}
\newfontfamily{\thaifont}{Norasi}[Script=Thai]
\newfontfamily{\thaifontsf}{Garuda}[Script=Thai]
\newfontfamily{\thaifonttt}{Noto Sans Thai}
\XeTeXlinebreaklocale "th"

\setotherlanguage{english}
\newfontfamily\englishfont{FreeSerif}[Ligatures=TeX]
\renewfontfamily\sffamilylatin{FreeSans}[Ligatures=TeX]
\renewfontfamily\ttfamilylatin{FreeMono}

\usepackage{hyperref}
\renewcommand{\UrlFont}{\ttfamilylatin}

\title{ตัวอย่างเช่นไทย}
\author{Lim Lian Tze}

\begin{document}

\maketitle

\tableofcontents

\section{บทนำ}
คอรัปชันจุ๊ยโปรดิวเซอร์ สถาปัตย์จ๊าบ แจ็กพ็อต ม้าหินอ่อน ซากุระคันถธุระ ฟีดสตาร์ท งี้ บอยคอตอิ่มแปร้สังโฆคำสาปแฟนซี ศิลปวัฒนธรรมไฟลท์จิ๊กโก๋กับดัก เจลพล็อตมาม่าซากุระดีลเลอร์ ซีนดัมพ์ แฮปปี้ เอ๊าะอุรังคธาตุซิม ฟินิกซ์เทรลเล่อร์อวอร์ด แคนยอนสมาพันธ์ ครัวซองฮัมอาข่าเอ็กซ์เพรส 

\begin{otherlanguage}{english}
Some maths and some English: $\sqrt{x + y - \alpha \times 2}$ 
\end{otherlanguage}

\textsf{ฟินิกซ์เทรลเล่อร์อวอร์ด} \textenglish{\sffamily `Sans-serif' Thai font!}

\begin{otherlanguage}{english}
A list of available Thai fonts on the Overleaf server can be found at \url{https://www.overleaf.com/help/193#!Thai}. You can upload your own .otf and .ttf fonts to use them with XeLaTeX.
\end{otherlanguage}

\end{document}